% from ega0-23-fr.tex
\oldpage[0\textsubscript{IV-1}]{220}

fractions $R$ est donc composé direct d'un nombre fini de corps $L_i$~;
l'image canonique $B_i$ de $B$ dans $L_i$ est un anneau local noethérien
\emph{intègre et complet}, dont $L_i$ est le corps de fractions, et la
fermeture intégrale $B'$ de $B$ dans $R$ est \emph{composée directe} des
$B'_i$, où $B'_i$ est la clôture intégrale de $B_i$. Mais en vertu de
(23.1.5), $B'_i$ est une $B_i$-algèbre \emph{finie}, donc un anneau local
\emph{noethérien} intégralement clos. On sait (Bourbaki,
\emph{Alg. comm.}, chap.~VII, §~1, n\textsuperscript{o}~3, cor. du th.~2) que
$B'_i$ est un anneau de Krull. Or, pour tout $i$, on a un homomorphisme
$\varphi_i:K\to L_i$, qui est injectif, et par suite
$\varphi_i^{-1}(B'_i)$ est un anneau de Krull dans $K$. Mais comme
$A'=B'\cap K$ en vertu du lemme et que
$B'=\bigcap_i\varphi_i^{-1}(B'_i)$, $A'$ est intersection d'un nombre fini
d'anneaux de Krull et est par suite un anneau de Krull.

Pour prouver la dernière assertion de la proposition, notons qu'il existe une
$A$-algèbre finie $C\subset K$ telle que le morphisme
$\operatorname{Spec}(A')\to\operatorname{Spec}(C)$ soit radiciel et un
homéomorphisme (23.2.4)~; pour tout idéal premier
$\mathfrak p'\in\operatorname{Spec}(A')$, $\mathfrak p'$ est donc le seul
idéal premier de $A'$ au-dessus de $\mathfrak p=\mathfrak p'\cap C$ et
$A'_{\mathfrak p'}$ la clôture intégrale de $C_{\mathfrak p}$~; d'autre part
le fait que $\operatorname{Spec}(A')\to\operatorname{Spec}(C)$ est un
homéomorphisme entraîne que l'application
$\mathfrak p'\mapsto\mathfrak p'\cap C$ est une bijection de l'ensemble des
idéaux premiers de hauteur $1$ de $A'$ sur l'ensemble des idéaux premiers de
hauteur $1$ de $C$~; d'où la conclusion, compte tenu de Bourbaki,
\emph{Alg. comm.}, chap.~VII, §~1, n\textsuperscript{o}~6, th.~4.

\begin{remarks}[23.2.8]
\label{0.23.2.8-fr}
\begin{enumerate}
  \item[(i)] On aura soin de noter, dans l'application de (23.2.6) et
  (23.2.7), que l'anneau $A'$ \emph{n'est pas nécessairement un anneau
  noethérien}, comme le montre un exemple de Nagata avec $\dim(A)=3$ [30].

  \item[(ii)] Les conclusions de (23.2.7) sont encore valables si $A'$ est
  la fermeture intégrale de $A$ dans une extension finie $K'$ de $K$~; il
  suffit en effet de considérer une $A$-algèbre finie $B$ dont $K'$ est le
  corps des fractions~; $B$ est un anneau semi-local noethérien, donc
  intersection d'un nombre fini d'anneaux locaux noethériens
  $B_{\mathfrak m_i}$ ($\mathfrak m_i$ idéaux maximaux de $B$), et sa clôture
  intégrale (qui est égale à $A'$) est intersection des
  $A'_{\mathfrak m_i}$ (Bourbaki, \emph{Alg. comm.}, chap.~II, §~3,
  n\textsuperscript{o}~3, cor.~4 du th.~1) qui sont les clôtures intégrales
  des $B_{\mathfrak m_i}$, donc des anneaux de Krull par (23.1.13)~; par
  suite $A'$ est un anneau de Krull.

  \item[(iii)] On peut montrer ([30], 3.3.10) que pour tout anneau noethérien
  intègre $A$ (non nécessairement local), la clôture intégrale de $A$ est un
  anneau de Krull.
\end{enumerate}
\end{remarks}

\begin{corollary}[23.2.9]
\label{0.23.2.9-fr}
Soient $A$ un anneau noethérien intègre, $A'$ sa clôture intégrale. Supposons
que pour tout anneau $C$ tel que $A\subset C\subset A'$ et tel que $C$ soit
une $A$-algèbre finie, et pour tout idéal premier $\mathfrak q$ de hauteur $1$
dans $C$, $\mathfrak q\cap A$ soit un idéal premier de hauteur $1$ dans $A$.
Soit $P$ l'ensemble des idéaux premiers de hauteur $1$ dans $A$~; on a alors
\[
  A'=\bigcap_{\mathfrak p\in P}A'_{\mathfrak p}.
\]
\end{corollary}

\begin{proof}
Comme $A'$ est un $A$-module sans torsion, on a
$A'=\bigcap A'_{\mathfrak m}$, où $\mathfrak m$ parcourt l'ensemble des idéaux
maximaux de $A$ (Bourbaki, \emph{Alg. comm.}, chap.~II, §~3,
n\textsuperscript{o}~3, cor.~4 du th.~1)~;

